Source : D. Carey et J. Rabesona, « Tax Ratios on Labour and Capital Income and on Consumption », OECD Economic Studies n° 35, 2002, tableau A2 (méthode révisée).\par Note : moyennes par période des pays dont les données sont complètes depuis 1975. Le ratio rapporte tous les impôts sur les revenus du capital, y compris ceux des ménages, au revenu du capital mesuré par la comptabilité nationale. Il dépend de conventions : si l'on attribue au travail une partie du revenu des indépendants, la hausse moyenne est réduite d'environ deux points, et celle de la France presque entièrement.
